% ID: SEL-D0082-OA-E04
% ESTADO: cotejada. Segunda pasada visual y borrador compilado revisados.
% Fuente DOC: Selectividad/sel_2012_matematicas_aplicadas/4 MATEMÁT. APLIC. CC. SS/EXAMEN 4-ANDALUCÍA.doc
% SHA-256 original: 8446c5da6b4b7cfc0f48b89d1b2d052ae0c8a1c4b5bc0b86b7e63399ce8656f8
% Página convertida: 1; opción A; ejercicio 4.

De una muestra aleatoria de 120 alumnos presentados a las Pruebas de Acceso, sólo 15 han resultado no aptos.
\begin{enumerate}[label=\alph*)]
\item \textbf{(1.5 puntos)} Calcule un intervalo de confianza, al $99\%$, para estimar la proporción de alumnos que han resultado aptos en dicha prueba.
\item \textbf{(1 punto)} Manteniendo la misma confianza, \textquestiondown{}cuál debe ser el tamaño mínimo de la muestra para estimar la proporción de alumnos aptos, cometiendo un error inferior al $5\%$?
\end{enumerate}
